\chapter{Proceso de desarrollo del proyecto}
\label{cap:metodologia}

Este capítulo describe cómo se ha construido el proyecto a lo largo del tiempo. Tras situar la organización del código, el grueso del capítulo ---y su verdadero contenido--- es el recorrido \textbf{sprint a sprint}: qué se hizo en cada etapa, qué tecnologías se eligieron frente a qué alternativas y qué lecciones dejó cada una. El balance cuantitativo del proyecto (horas dedicadas, planificación temporal, presupuesto y riesgos) se recoge en el Capítulo~\ref{cap:conclusiones}, junto al resto de conclusiones.

El desarrollo se ha conducido siguiendo un marco de trabajo ágil (\textit{Agile}), caracterizado por entregas iterativas e incrementales. Dada la naturaleza cambiante del ámbito de la Inteligencia Artificial y de los Modelos de Lenguaje Grande (LLMs), la flexibilidad para revisar decisiones técnicas, sustituir herramientas o ajustar el encaminamiento resultó especialmente importante durante el desarrollo.

Se estableció un ciclo de vida con las siguientes fases:
\begin{enumerate}
    \item \textbf{Elicitación y Análisis de Requisitos:} Identificación de fuentes de información clave (SharePoint, páginas web corporativas y el BOE), así como las restricciones de privacidad requeridas por un despliegue puramente \textit{on-premise}.
    \item \textbf{Diseño de Arquitectura:} Planificación de la infraestructura basada en microservicios contenerizados y selección del stack tecnológico idóneo.
    \item \textbf{Implementación Incremental:} Programación progresiva de los módulos de ingesta, OCR, embeddings, búsqueda de Qdrant, backend y, por último, de la interfaz de usuario de OpenWebUI.
    \item \textbf{Pruebas y Validación:} Validación funcional del encaminamiento inteligente y de las capacidades del RAG mediante métricas comparativas empíricas y experimentación cualitativa.
\end{enumerate}

El trabajo se dividió en \textbf{siete sprints} de aproximadamente un mes de duración, desarrollados entre octubre de 2025 y mayo de 2026, a los que siguió una fase final de documentación, redacción de la memoria y difusión del proyecto hasta julio de 2026. Cada sprint se centró en un módulo funcional del sistema y terminó con un entregable operativo. La distribución temporal completa (diagrama de Gantt) y las horas dedicadas a cada sprint se presentan en la Sección~\ref{sec:esfuerzo}.

% ============================================================
\section{Organización del Código}
\label{sec:organizacion-codigo}
% ============================================================

Antes de describir cada sprint conviene tener presente la estructura del repositorio del proyecto, porque cada directorio se corresponde de forma bastante directa con el sprint en el que nació. Esta correspondencia entre organización del código y evolución temporal facilita tanto el mantenimiento como la lectura del historial del proyecto:

\begin{lstlisting}[numbers=none, basicstyle=\ttfamily\footnotesize, caption={Estructura del repositorio y sprint en el que se desarrolló cada componente.}, label={lst:estructura-repo}]
TFG-RAG-URJC/
|-- backend/          <- Servicio FastAPI (sprints 2, 4, 6 y 7)
|   `-- app/          <- api, core (RAG, SQL, QueryProcessor),
|                        integrations, processing, storage...
|-- services/
|   |-- ollama/       <- Motor de inferencia local (sprint 1)
|   |-- litellm/      <- Proxy unificado de modelos (sprint 1)
|   |-- indexer/      <- Sincronizacion SharePoint+OCR (sprint 3)
|   `-- openwebui/    <- Interfaz y pipeline JARVIS (sprint 4)
|-- mcp-boe-server/   <- Servidor MCP del BOE (sprint 5)
|-- scripts/          <- Fine-tuning, datasets, backups (spr. 5)
|-- config/           <- NGINX, Prometheus, Grafana (spr. 1 y 5)
|-- database/         <- Esquema inicial PostgreSQL (sprint 1)
|-- tests/            <- Pruebas del sistema (transversal)
|-- docs/             <- Documentacion tecnica (transversal)
|-- github_pages/     <- Pagina web del proyecto
`-- memoria/          <- Esta memoria (LaTeX)
\end{lstlisting}

% ============================================================
\section{Los Sprints en Detalle}
\label{sec:sprints}
% ============================================================

A continuación se define cada sprint siguiendo un esquema común: descripción y objetivo, tareas realizadas, tecnologías elegidas frente a las alternativas valoradas, y estado final con las lecciones aprendidas. Todos los sprints finalizaron con su entregable operativo, por lo que el estado se detalla junto a las lecciones de cada uno.

\subsection{Sprint 1 (octubre de 2025): Infraestructura base}

\textbf{Descripción.} Construir los cimientos del sistema: la plataforma de contenedores, los servicios de datos y el motor de inferencia local, validando que era viable ejecutar un LLM útil en el hardware disponible.

\textbf{Tareas.} Definición del \texttt{docker-compose} inicial y de la red interna aislada; despliegue de PostgreSQL con su esquema inicial (\texttt{database/init.sql}), Qdrant y Redis; instalación de Ollama con los primeros modelos (Llama 3.1) y del proxy LiteLLM; pruebas de concepto de chat local.

\textbf{Tecnologías y alternativas valoradas.}
\begin{itemize}
    \item \textbf{Ollama} frente a \texttt{llama.cpp} puro o vLLM: Ollama aporta gestión de modelos, carga/descarga automática de VRAM y una API estable con mínima configuración. vLLM ofrece mayor rendimiento bruto, pero está orientado a servir un único modelo de forma intensiva; aquí se necesitaba rotar varios modelos sobre una sola GPU.
    \item \textbf{Qdrant} frente a pgvector, Chroma o Pinecone: se buscaba una base vectorial \textit{on-premise} con índice HNSW, colecciones separadas y filtrado por metadatos (la base del control de acceso por departamento). Pinecone se descartó por ser un servicio en la nube, contrario al requisito de soberanía del dato.
    \item \textbf{Docker Compose} frente a Kubernetes: para un despliegue de nodo único, Kubernetes añadía complejidad operativa sin beneficio real.
\end{itemize}

\textbf{Estado y lecciones.} Completado. La lección central llegó ya en este sprint: la GPU es un recurso compartido y escaso, y esa restricción condicionó todas las decisiones posteriores (cuantización, modelos de embeddings en CPU, carga bajo demanda).

\subsection{Sprint 2 (noviembre de 2025): Backend y pipeline RAG}

\textbf{Descripción.} Desarrollar el motor del sistema: un backend propio con el pipeline RAG completo, desde la ingesta del documento hasta la respuesta con citas.

\textbf{Tareas.} API REST con FastAPI; ingesta de PDFs con \texttt{pdfplumber}; \textit{chunking} con solapamiento (500 tokens, solapamiento de 50); generación de embeddings con \path{paraphrase-multilingual-MiniLM-L12-v2} en CPU; búsqueda semántica en Qdrant; construcción del prompt restrictivo con citas de fuente y página.

\textbf{Tecnologías y alternativas valoradas.}
\begin{itemize}
    \item \textbf{FastAPI} frente a Flask o Django: asincronía nativa (esencial para no bloquear durante inferencias largas), validación con Pydantic y documentación OpenAPI generada automáticamente.
    \item \textbf{Pipeline RAG propio} frente a LangChain o LlamaIndex: los \textit{frameworks} aceleran el prototipo, pero ocultan los pasos intermedios. Implementar el pipeline a mano dio control fino sobre chunking, búsqueda y prompt, y sobre todo trazabilidad completa para depurar; esa comprensión del flujo es además parte del valor académico del trabajo.
    \item \textbf{Embeddings en CPU} frente a GPU: el modelo MiniLM es lo bastante ligero para ejecutarse en CPU, reservando la VRAM para los modelos generativos.
\end{itemize}

\textbf{Estado y lecciones.} Completado. La calidad del RAG depende tanto de la ingesta (tamaño de fragmento, solapamiento, limpieza) como del propio modelo generativo: gran parte de las mejoras posteriores vinieron de recuperar mejor, no de generar mejor.

\subsection{Sprint 3 (diciembre de 2025): SharePoint y OCR}

\textbf{Descripción.} Conectar el sistema a las fuentes documentales corporativas reales: la biblioteca SharePoint de la organización y los documentos escaneados.

\textbf{Tareas.} Servicio Indexer con sincronización periódica (cron cada 5 minutos); autenticación MSAL con flujo de credenciales de cliente y descarga incremental vía Microsoft Graph API; motor OCR PaddleOCR con aceleración GPU y \textit{fallback} automático a CPU; paralelización del OCR con Ray.

\textbf{Tecnologías y alternativas valoradas.}
\begin{itemize}
    \item \textbf{PaddleOCR} frente a Tesseract: mejor precisión sobre los documentos escaneados reales del entorno de validación y soporte GPU nativo.
    \item \textbf{Sondeo incremental} frente a \textit{webhooks} de Graph: los webhooks exigen exponer un endpoint público accesible desde Microsoft, incompatible con un despliegue cerrado; el sondeo cada pocos minutos resultó suficiente para la frecuencia de cambio real de los documentos.
\end{itemize}

\textbf{Estado y lecciones.} Completado. Al paralelizar la ingesta apareció una \textit{race condition} en la generación de embeddings que hubo que serializar; y quedó claro que las fuentes corporativas cambian constantemente, lo que convirtió la sincronización incremental en pieza central del diseño.

\subsection{Sprint 4 (enero de 2026): Agente JARVIS, scraping y búsqueda web}

\textbf{Descripción.} Dar al sistema su cara visible y sus capacidades web: el pipeline JARVIS sobre OpenWebUI, que enruta cada consulta al modo adecuado, y la extracción de contenido de internet.

\textbf{Tareas.} Desarrollo del pipeline JARVIS como \textit{hook} de OpenWebUI con detección de intenciones por reglas y sistema de \textit{valves} de configuración; scraping con Playwright (renderizado JavaScript completo) y cadena de extracción Trafilatura $\rightarrow$ Readability $\rightarrow$ BeautifulSoup; búsqueda en internet con DuckDuckGo.

\textbf{Tecnologías y alternativas valoradas.}
\begin{itemize}
    \item \textbf{OpenWebUI con pipeline propio} frente a interfaz web desarrollada desde cero: OpenWebUI ya resuelve autenticación, historial, streaming y gestión de usuarios; el valor diferencial estaba en la lógica de enrutamiento, no en reconstruir un chat.
    \item \textbf{Playwright} frente a peticiones HTTP con BeautifulSoup: gran parte de las páginas corporativas y de noticias actuales requieren ejecutar JavaScript para mostrar su contenido.
    \item \textbf{Router por reglas} frente a un LLM decisor: las reglas son instantáneas, gratuitas y deterministas, y el modelo local desplegado no ofrecía \textit{function calling} fiable. Esta decisión se analiza en detalle en el Capítulo~\ref{cap:implementacion}.
\end{itemize}

\textbf{Estado y lecciones.} Completado. Dos incidencias de despliegue dejaron lección: un pipeline duplicado por el mecanismo de descubrimiento de OpenWebUI y un script de inicialización de Ollama roto por los finales de línea CRLF al editar desde Windows; ambas reforzaron la necesidad de probar el despliegue completo desde cero, no solo el código.

\subsection{Sprint 5 (febrero de 2026): BOE, \textit{fine-tuning} y observabilidad}

\textbf{Descripción.} Ampliar el sistema en tres frentes: la fuente normativa oficial (BOE), la adaptación del modelo al dominio mediante \textit{fine-tuning} y la monitorización del conjunto.

\textbf{Tareas.} Servidor MCP del BOE e integración funcional por API; generación de un dataset sintético de preguntas y respuestas a partir de los documentos indexados en Qdrant; entrenamiento de un adaptador LoRA (PEFT, bitsandbytes, TRL) que dio lugar a \texttt{rag-qwen-ft}, el modelo JARVIS; despliegue del stack de observabilidad (Prometheus, Grafana, exportador DCGM para GPU, y posteriormente Loki y Promtail para logs).

\textbf{Tecnologías y alternativas valoradas.}
\begin{itemize}
    \item \textbf{LoRA} frente a \textit{fine-tuning} completo: ajustar todos los pesos de un modelo de 7B parámetros excede la VRAM disponible; LoRA entrena una fracción mínima de parámetros con resultados comparables para adaptación de dominio.
    \item \textbf{MCP} frente a integración directa por API: se implementaron ambas. MCP se adoptó como apuesta de estandarización futura, pero la interfaz web principal no lo soporta de forma nativa, por lo que la vía de producción siguió siendo la API (la valoración completa está en el Anexo~\ref{anexo:integraciones}).
    \item \textbf{Prometheus y Grafana} frente a soluciones de monitorización en la nube: coherencia con el requisito \textit{on-premise} del proyecto.
\end{itemize}

\textbf{Estado y lecciones.} Completado. La lección más importante de todo el proyecto salió de aquí: el \textit{fine-tuning} rindió resultados modestos para RAG, mientras que las mejoras de recuperación (búsqueda híbrida y reranking) tuvieron un impacto muy superior. Esto reorientó el esfuerzo restante hacia la calidad de la recuperación.

\subsection{Sprint 6 (marzo--abril de 2026): Agente SQL, seguridad y v2.2.0}

\textbf{Descripción.} Incorporar la consulta de datos estructurados en lenguaje natural y endurecer la seguridad del sistema, además de corregir defectos detectados en producción.

\textbf{Tareas.} Agente NL$\rightarrow$SQL sobre PostgreSQL con lista blanca de operaciones (\texttt{SELECT}) y auto-corrección ante errores; validación de tokens JWT de Azure AD contra JWKS en el backend; diagnóstico y corrección de la expansión de consultas asíncrona (versión 2.2.0).

\textbf{Tecnologías y alternativas valoradas.}
\begin{itemize}
    \item \textbf{Validación por lista blanca y expresiones regulares} frente a confiar en las instrucciones del prompt: un LLM puede ser inducido a generar SQL destructivo; la validación programática actúa como defensa en profundidad independiente del modelo.
    \item \textbf{Validación JWT en el backend} frente a delegar la autenticación en la interfaz: cada microservicio verifica la identidad por sí mismo, de modo que un fallo en la capa de presentación no expone los datos.
\end{itemize}

\textbf{Estado y lecciones.} Completado. La incidencia más formativa: la expansión de consultas llevaba tiempo desactivada por un defecto de configuración sin que ninguna prueba lo detectara. La lección es que los tests deben cubrir también la configuración efectiva del sistema en ejecución, no solo la lógica del código.

\subsection{Sprint 7 (mayo de 2026): QueryProcessor y v2.3.0}

\textbf{Descripción.} Cerrar el sistema con la capa de orquestación inteligente: decidir, para cada consulta, qué modelo y qué estrategia de recuperación usar.

\textbf{Tareas.} Módulo \texttt{QueryProcessor} con detección de intención, encaminamiento inteligente entre modelos (JARVIS para consultas factuales y procedimentales; Qwen 2.5 32B para las analíticas) y expansión de consultas multivariante; summarización del historial de conversación con el modelo de 32B en el \texttt{MemoryManager} (versión 2.3.0).

\textbf{Tecnologías y alternativas valoradas.}
\begin{itemize}
    \item \textbf{Encaminamiento por intención entre dos modelos} frente a usar siempre el modelo mayor: las pruebas comparativas mostraron que el modelo ajustado de 7B domina en recuperación precisa con citas mientras que el de 32B destaca en razonamiento analítico; usar cada uno donde sobresale mejora calidad y latencia a la vez.
\end{itemize}

\textbf{Estado y lecciones.} Completado. El modelo más grande no siempre es la mejor opción: adaptar la elección del modelo a la intención de la consulta resultó más eficiente que cualquier elección fija.
